\section{Linear costs of root paths}
\label{sec:linear-paths}

The global constraint on a test has a concrete dual meaning.  A vector
can be paid for either as a Hilbert vector or as a sum of root paths;
the norm is the least sum of these two costs.  We first prove this exact
identification.  We then use its two descriptions in different ways:
clipping controls dual tests, whereas removal and prediction control
signed path coefficients.

\subsection{The norm and its finite heads}

For an integer $n\ge1$, let
$T_n=\bigcup_{j=0}^n\mathbb N^j$, with root $o=\varnothing$ and
prefix order $\preceq$.  The root belongs to every root path.  On the
real Hilbert space $H_n=\ell_2(T_n)$ put
\[
 b_s=\sum_{r\preceq s}e_r,
 \qquad
 B_n\mu=\sum_{s\in T_n}\mu_s b_s
 \quad(\mu\in\ell_1(T_n)).
\]
Since $\|b_s\|_2\le\sqrt{n+1}$, the series defining $B_n\mu$
converges absolutely in $H_n$ and
$\|B_n\|\le\sqrt{n+1}$.  Define
\begin{align}
 L_n(x)&=\inf_{x=h+B_n\mu}\bigl(\|h\|_2+\|\mu\|_1\bigr),
 \label{eq:linear-cost}\\
 p_n(f)&=\max\left\{\|f\|_2,
       \sup_{s\in T_n}\left|\sum_{r\preceq s}f_r\right|\right\}.
 \label{eq:linear-polar}
\end{align}
The infimum ranges over $h\in H_n$ and $\mu\in\ell_1(T_n)$;
attainment is not assumed.  Restricting $\mu$ to finite support gives
the same infimum: truncate an arbitrary $\mu$ in $\ell_1$ and absorb
the vanishing remainder $B_n(\mu-\mu^{(k)})$ into $h$.

\begin{proposition}[Exact duality for the linear cost]
\label{prop:linear-duality}
The space $X_n=(H_n,L_n)$ is complete and
\[
 \frac{\|x\|_2}{\sqrt{n+1}}\le L_n(x)\le\|x\|_2,
 \qquad
 L_n(x)=\sup_{p_n(f)\le1}\langle f,x\rangle.
\]
Every coordinate satisfies $|x_s|\le L_n(x)$.
Its continuous dual is $H_n$, under the Hilbert pairing, with norm
$p_n$.  In particular $L_n$ is exactly the global test norm on $T_n$.
\end{proposition}

\begin{proof}
The choice $h=x$, $\mu=0$ proves the upper bound.  Every representation
of $x$ satisfies
\[
 \|x\|_2\le\|h\|_2+\sqrt{n+1}\|\mu\|_1
 \le\sqrt{n+1}(\|h\|_2+\|\mu\|_1),
\]
which proves the lower bound.  Scaling and adding representations,
chosen within an arbitrary error of their infima, prove homogeneity
and the triangle inequality.  The lower bound gives definiteness;
norm equivalence gives completeness and density of finite arrays.

For $p_n(f)\le1$, absolute convergence allows us to write
\[
 |\langle f,x\rangle|
 \le\|h\|_2+
       \sum_s|\mu_s|\left|\sum_{r\preceq s}f_r\right|
 \le\|h\|_2+\|\mu\|_1.
\]
Thus the displayed supremum is at most $L_n(x)$.  Conversely,
Hahn--Banach gives, for $x\ne0$, a functional $\varphi$ such that
$\varphi(x)=L_n(x)$ and $|\varphi(z)|\le L_n(z)$.  Since
$L_n(z)\le\|z\|_2$, it has the form $\varphi(z)=\langle f,z\rangle$
with $\|f\|_2\le1$.  The representation $b_s=B_ne_s$ gives
$L_n(b_s)\le1$, hence
$|\sum_{r\preceq s}f_r|=|\varphi(b_s)|\le1$.  This proves equality,
including the trivial case $x=0$.

For an arbitrary $f\in H_n$, the representation estimate gives dual
norm at most $p_n(f)$.  The Hilbert unit ball and every $\pm b_s$ belong
to the $L_n$ unit ball.  Testing on these vectors gives both reverse
bounds in \eqref{eq:linear-polar}, and therefore the exact dual norm.
\end{proof}

A set $A\subseteq T_n$ is an \emph{initial set} if it contains every
prefix of each of its vertices; the empty set is allowed.  Write $P_A$
for coordinate restriction to $A$.  Restricting a root path to $A$
gives either zero or the root path ending at its last vertex in $A$.
Grouping coefficients by this endpoint cannot increase their
$\ell_1$ norm.  Therefore $P_A$ is contractive on $X_n$.  It is also
contractive on its dual: restriction freezes every prefix sum at its
last retained vertex and decreases the Hilbert norm.

For a finite initial set $A$, the unit ball of the induced norm on
$\mathbb R^A$ is
\begin{equation}
 \operatorname{conv}\left(B_{\ell_2(A)}\cup
                         \{b_s,-b_s:s\in A\}\right).
 \label{eq:linear-finite-ball}
\end{equation}
Indeed this is a compact convex symmetric set whose polar is precisely
the set of tests in \eqref{eq:linear-polar} supported on $A$.
Finite-dimensional separation and Proposition~\ref{prop:linear-duality}
identify the balls.  The assertion of compactness concerns the finite
head only.  On the whole block the ball is the norm-closed convex hull
of the Hilbert ball and the signed paths: one inclusion follows from
convexity, and an almost optimal representation of a vector in the unit
ball, followed by finite truncation and a scaling tending to one,
proves the other.

At height one, this closed convex hull is exactly the unit ball used by
Dilworth et al.: the paths are the root singleton $e_o$, already in the
Hilbert ball, and the vectors $e_o+e_s$ for children $s$. Under the
identification of $o$ with their first coordinate, $p_1$ is also their
dual formula \cite[Lemma~2.5]{DKRRZ2016}. The higher blocks retain this
path geometry through arbitrarily many generations.

We use the outer Hilbert sum
\begin{equation}
 X=\left(\bigoplus_{n\ge1}X_n\right)_{\ell_2},
 \qquad \|x\|^2=\sum_nL_n(x_n)^2.
 \label{eq:linear-outer-sum}
\end{equation}
A finite initial subset of the disjoint union of the $T_n$ will be
called a \emph{finite head}.  Its restriction $P_A$ is contractive on
$X$, and such restrictions converge strongly to the identity along an
exhaustion.  To see the latter assertion, first truncate the block sum,
then approximate the remaining coordinates by finite arrays and enlarge
their support to an initial set.  If $A\subseteq C$ are initial sets,
then the window projection
\[
 R_{C\setminus A}=P_C-P_A
 \quad\hbox{satisfies}\quad \|R_{C\setminus A}\|\le2.
\]
In particular the complementary projection $I-P_A$ has norm at most
two; its contractivity is not needed.

\begin{proposition}[Duality and branching in the outer sum]
\label{prop:linear-geometry}
The space $X$ is reflexive, has separable dual, and its dual is the
Hilbert sum of the spaces $(H_n,p_n)$, with pairing
\[
 f(x)=\sum_n\langle f_n,x_n\rangle,
 \qquad \|f\|_*^2=\sum_np_n(f_n)^2.
\]
Every coordinate vector $e_{n,s}$ and every root path $b_{n,s}$ has norm
one.  For fixed $n$ and $|s|<n$, the children
$(e_{n,s^\frown j})_{j\ge1}$ are weakly null.
If $\alpha\|\cdot\|\le N\le\beta\|\cdot\|$, with
$0<\alpha\le\beta$, then $\overline\delta_N(t)=0$ for every
$0<t\le\alpha/\beta$.  In particular no equivalent norm on $X$
is asymptotically uniformly convex.
\end{proposition}

\begin{proof}
Each block is an equivalently normed Hilbert space and hence reflexive.
For completeness of the sum, a Cauchy sequence has a limit in each
block; its norm bounds and Cauchy estimates pass to the limit first on
finite sets of blocks and then by increasing those sets.  Cauchy--Schwarz
shows that each displayed dual family defines a bounded functional.
Conversely, restrict an arbitrary functional to finitely many blocks,
choose almost norming vectors there, and weight them proportionally to
the block dual norms.  This proves that the sum of their squares is
bounded by the square of the original functional norm.  Density of
finite-block vectors proves the representation and equality of norms.
Applying the same argument to the dual sum proves surjectivity of the
canonical map into the bidual.  Finite arrays are dense in the dual
by the same block truncation and Hilbert equivalence, so it is separable.

The cost formula gives $L_n(e_s),L_n(b_s)\le1$.  The dual vector $e_s$
has $p_n(e_s)=1$ and evaluates to one on both vectors, giving equality.
The weak nullity assertion follows from square summability of every
fixed dual array on the $n$th block.

For clarity we give the renorming argument with its scale and gain.
Suppose $\alpha\|\cdot\|\le N\le\beta\|\cdot\|$, where
$0<\alpha\le\beta$, and suppose $d=\overline\delta_N(t)>0$ for
some $0<t\le\alpha/\beta$.  Choose $0<\eta<d/2$.
For each unit center $u$ there is a finite-codimensional closed
subspace $F$ such that $N(u+tv)\ge1+2\eta$ on its unit sphere.
Convexity on a ray gives $N(u+v)\ge1+2\eta$ whenever $v\in F$
and $N(v)\ge t$.

At an internal vertex $s$ take
$u=b_{n,s}/N(b_{n,s})$ and
$z_j=e_{n,s^\frown j}/N(b_{n,s})$.  These vectors are weakly null and
$N(z_j)\ge\alpha/\beta\ge t$.  Their quotient images in $X/F$
tend to zero, since that quotient is finite-dimensional.  Choose
$v_j\in F$ with $N(v_j-z_j)\to0$.  If $N(v_j)<t$, replace $v_j$
by $tv_j/N(v_j)$; for large $j$ this is defined, and the extra change
is at most $N(v_j-z_j)$.  Thus $z_j$ is approximated by vectors in
$F$ of norm at least $t$.  For a sufficiently large child index,
\[
 N(b_{n,s^\frown j})\ge(1+\eta)N(b_{n,s}).
\]
Iterating to depth $n$ gives $\alpha(1+\eta)^n\le\beta$, impossible
for all $n$.  The asymptotic modulus is nonnegative: the kernel of a
norming functional at a unit center has $N(u+tv)\ge1$.
This proves the claimed zero value at every indicated scale.  Taking
$\eta=d/4$ gives the often useful growth factor $1+d/4$.
\end{proof}

Two distinctions about the model will matter below.  The global test
norm on the infinite-height tree restricts to $L_n$ on $T_n$: restrict
an admissible test in one direction, and extend by zero in the other,
which freezes all longer prefix sums.  This identifies the completed
finite-height block, but does not identify the infinite-height
completion with \eqref{eq:linear-outer-sum}.  Also the quadratic cost
\[
 Q_n(x)=\inf_{x=h+B_n\mu}
           \bigl(\|h\|_2^2+\|\mu\|_1^2\bigr)^{1/2}
\]
satisfies $Q_n\le L_n\le\sqrt2Q_n$ and is different from $L_n$:
$L_n(e_o)=1$, whereas $Q_n(e_o)=1/\sqrt2$.  The upper bound in the
last equality comes from splitting the root coordinate equally; the
lower bound follows from
$1\le\|h\|_2+\|\mu\|_1\le\sqrt2(\|h\|_2^2+\|\mu\|_1^2)^{1/2}$.
Thus constants proved for one cost cannot be transferred by identifying
the two norms.

\subsection{Clipping the tail of a dual test}
\label{subsec:linear-clipping}

A head test has a constant prefix sum after leaving the head. We split
the coefficients of a second test outside the head into two pieces whose prefix sums point towards zero
relative to this constant.  Hilbert orthogonality pays for their
coefficients, and the inward direction pays for the prefix constraint.

\begin{lemma}[Inward decomposition for the linear polar]
\label{lem:linear-inward}
Let $A$ be a finite head.  Suppose $h,g\in X^*$ have norms at most one
and $P_Ah=h$.  There are $r_1,r_2\in X^*$, vanishing on $A$, such that
\[
 r_1-r_2=(I-P_A)g,\qquad \|r_l\|_*\le2,
 \qquad
 \|\lambda h+\mu r_l\|_*^2\le\lambda^2+4\mu^2
\]
for $l=1,2$ and $\lambda,\mu\ge0$.
\end{lemma}

\begin{proof}
In block $n$ write $A_n=A\cap T_n$ and set
\[
 H(s)=\sum_{r\preceq s}h_{n,r},\qquad
 V(s)=\sum_{\substack{r\preceq s\\r\notin A_n}}g_{n,r},
 \qquad a_n=p_n(h_n),\quad c_n=p_n(g_n).
\]
Here $|H|\le a_n$.  The sum $V$ is a difference of two prefix sums
of $g_n$, or a single prefix sum if the path misses $A_n$, so
$|V|\le2c_n$.  It vanishes on $A_n$.  Define
\[
 (V_1,V_2)=
 \begin{cases}
  (\min(V,0),-\max(V,0)),&H\ge0,\\
  (\max(V,0),-\min(V,0)),&H<0.
 \end{cases}
\]
Let $r_{l,n,s}=V_l(s)-V_l(s^-)$, with predecessor value zero at the
root.  These coefficients vanish on $A_n$.  Outside $A_n$, the head
sum $H$ is constant along each remaining branch, and both scalar
clipping maps are $1$-Lipschitz.  Therefore
\[
 |r_{l,n,s}|\le|g_{n,s}| \quad(s\notin A_n).
\]
At a first exit vertex this follows from the zero predecessor value
of $V_l$.  Moreover the prefix sums of $r_{l,n}$ are $V_l$, and
$V_1-V_2=V$, proving the claimed difference.

The supports of $h_n$ and $r_{l,n}$ are disjoint.  Hence their Hilbert
norm satisfies
\[
 \|\lambda h_n+\mu r_{l,n}\|_2^2
 \le\lambda^2a_n^2+\mu^2c_n^2.
\]
The products $HV_l$ are nonpositive, so each prefix sum satisfies
\[
 |\lambda H+\mu V_l|
 \le\max\{\lambda a_n,2\mu c_n\}.
\]
These two estimates give
$p_n(\lambda h_n+\mu r_{l,n})^2
 \le\lambda^2a_n^2+4\mu^2c_n^2$.
Summing over $n$ proves the asserted inequality; taking
$\lambda=0$, $\mu=1$ proves $r_l\in X^*$ and its norm bound.
\end{proof}

\begin{theorem}[Clipping estimate on a head and its complement]
\label{thm:linear-head-tail-clipping}
If $x,y\in X$, $P_Ax=x$ and $P_Ay=0$ for a finite head $A$, then,
for every $c>0$,
\begin{equation}
 \max\{\|x+y\|,\|x-y\|\}
 \ge\frac{\|x\|+c\|y\|/2}{\sqrt{1+4c^2}}.
 \label{eq:linear-clipping-family}
\end{equation}
In particular the left side is at least
$\sqrt{\|x\|^2+\|y\|^2/16}$.
\end{theorem}

\begin{proof}
Choose norming dual vectors $h,g$ for $x,y$, with norms at most one,
and replace $h$ by $P_Ah$.  The projection preserves its evaluation
at $x$ and does not increase its norm.  Lemma~\ref{lem:linear-inward}
gives $r_1(y)-r_2(y)=\|y\|$, so at least one of
$|r_1(y)|,|r_2(y)|$ is at least $\|y\|/2$.  For this index choose
a sign $\varepsilon$ with $\varepsilon r_l(y)\ge\|y\|/2$.
Because $h(y)=r_l(x)=0$, evaluation of $h+cr_l$ on
$x+\varepsilon y$ proves \eqref{eq:linear-clipping-family}.
Its supremum over $c>0$ is the displayed square root, by
Cauchy--Schwarz in two dimensions; endpoint cases follow by a limit.
\end{proof}

The next result has an arbitrary center and no chosen head in its
statement.  Weak nullity removes the head coordinates of the moving
vectors; the approximation of the center must be made before taking
the lower limit.

\begin{theorem}[Clipping along a weakly null sequence]
\label{thm:linear-weak-sequence}
Let $x\in X$ and let $(y_j)$ be a bounded weakly null sequence in $X$.
If $L=\liminf_j\|y_j\|$, then
\[
 \liminf_j\max\{\|x+y_j\|,\|x-y_j\|\}
 \ge\sqrt{\|x\|^2+L^2/16}.
\]
\end{theorem}

\begin{proof}
Fix $\eta>0$ and choose a finite head $A$ with
$\|x-P_Ax\|\le\eta$, using the strong head exhaustion. Put
$z_j=(I-P_A)y_j$. Since $P_A$ has finite rank and $y_j$ is weakly null,
$\|P_Ay_j\|\to0$. In particular,
\[
 \bigl|\|z_j\|-\|y_j\|\bigr|\le\|P_Ay_j\|\longrightarrow0.
\]
Apply Theorem~\ref{thm:linear-head-tail-clipping} to $P_Ax$ and $z_j$.
The triangle inequality then gives
\[
 \max\{\|x+y_j\|,\|x-y_j\|\}
 \ge\sqrt{\|P_Ax\|^2+\|z_j\|^2/16}
       -\eta-\|P_Ay_j\|.
\]
Keep this head fixed while taking the lower limit in $j$. The result is
at least $\sqrt{\|P_Ax\|^2+L^2/16}-\eta$. Finally let
$\eta\downarrow0$, choosing the heads along the strong exhaustion, so
that $\|P_Ax\|\to\|x\|$. This proves the stated bound.
\end{proof}

\subsection{Path coefficients on complementary forests}
\label{sec:linear-positive-mass}

The clipping estimate already controls a displacement supported outside
a head. We now examine the path representations in
\eqref{eq:linear-cost} themselves. The deficit between representation
cost and the value of a head norming functional on the represented vector
bounds the mass whose sign opposes that
functional. The remaining endpoint coefficients in each complementary
tree have one sign. Positive-mass monotonicity
uses that sign directly; the next subsection removes a controlled amount
of mass so that the remainder is small in Hilbert norm. These arguments
record their own tail bounds and expose coefficient mechanisms that will
reappear in the variation and stopping arguments.

Fix one tree \(T_n\), a head \(A\), and \(U=T_n\setminus A\).
The set \(U\) is a forest whose roots are the first vertices outside
\(A\). Define its intrinsic path vectors and norm by
\begin{align*}
 b_t^U&=\sum_{v\in U,\ v\preceq t}e_v,
 &B_U\alpha&=\sum_{t\in U}\alpha(t)b_t^U,\\
 L_U(z)&=\inf_{z=h+B_U\alpha}(\|h\|_2+\|\alpha\|_1).
\end{align*}
All paths have at most \(n+1\) vertices, so these expressions define a
bounded map from \(\ell_1(U)\) to \(\ell_2(U)\) and a Hilbert-equivalent
norm. If \(v\) is supported in \(U\), then
\begin{equation}\label{eq:linear-forest-comparison}
 L_n(v)\le2L_U(v|_U).
\end{equation}
Indeed, after extension by zero, \(b_t^U=b_t-P_A b_t\), of cost at most
two, while a Hilbert residual extends without changing its cost. Changing
sign on any collection of whole trees of \(U\) is an isometry for \(L_U\).

\begin{lemma}[Monotonicity of positive path mass]
\label{lem:linear-positive-monotonicity}
If \(0\le\nu\le\mu\) in \(\ell_1(U)\), then
\(L_U(B_U\nu)\le L_U(B_U\mu)\).
\end{lemma}
\begin{proof}
Choose a norming functional \(\phi\in\ell_2(U)\) for \(B_U\nu\).
Its dual norm is at most one, hence
\(\|\phi\|_2\le1\) and
\(F_t:=\sum_{v\in U,\,v\preceq t}\phi(v)\) satisfies \(|F_t|\le1\).
At a root set the value at its nonexistent parent equal to zero, and put
\[
 \psi(v)=|F_v|-|F_{\operatorname{par}(v)}|.
\]
Then \(|\psi(v)|\le|\phi(v)|\), and the prefix sums of \(\psi\) are
\(|F_t|\). Thus \(\psi\) also has dual norm at most one: pairing it
with a decomposition \(h+B_U\alpha\) costs at most
\(\|h\|_2+\|\alpha\|_1\). Absolute convergence of the coefficient
sums gives
\[
 L_U(B_U\nu)=\sum_t\nu(t)F_t
 \le\sum_t\mu(t)|F_t|
 =\langle B_U\mu,\psi\rangle\le L_U(B_U\mu).
\]
The claim concerns positive path coefficients; it asserts no general
coordinatewise monotonicity of the norm.
\end{proof}

\begin{theorem}[A tail estimate from positive mass]
\label{thm:linear-positive-tail}
Let \(C=6(\sqrt2+1)\). In one component, if
\(P_Ax=x\), \(P_Ay=0\), and \(L_n(x\pm y)\le1\), then
\begin{equation}\label{eq:linear-positive-unit}
 L_n(y)\le C\sqrt{1-L_n(x)}.
\end{equation}
For arbitrary \(x,y\) satisfying \(P_Ax=x\) and \(P_Ay=0\), writing
\(\beta=\max\{L_n(x+y),L_n(x-y)\}\) and \(r=L_n(x)\), one has
\begin{equation}\label{eq:linear-positive-scaled}
 L_n(y)^2\le C^2\beta(\beta-r)\le C^2(\beta^2-r^2).
\end{equation}
Consequently, for a finite head projection \(P\) on \(X\),
\begin{equation}\label{eq:linear-positive-outer}
 \|y\|^2\le C^2\bigl(\|x+y\|^2+\|x-y\|^2-2\|x\|^2\bigr)
 \quad(Px=x,\ Py=0).
\end{equation}
\end{theorem}
\begin{proof}
For the unit estimate set \(r=L_n(x)\) and choose a norming functional
\(g=gP_A\), of norm at most one, with \(g(x)=r\); use zero if \(x=0\).
For each sign choose
\(x\pm y=h_\pm+B_n\alpha_\pm\) of cost at most \(1+\eta\), where
\(\eta>0\). Put \(d=1+\eta-r\). Subtracting the value of \(g\) from
the cost gives nonnegative terms with sum at most \(d\):
\[
 \|h_\pm\|_2-g(h_\pm)
 +\sum_t\bigl(|\alpha_\pm(t)|-\alpha_\pm(t)g(b_t)\bigr)\le d.
\]
In particular,
\[
 \|h_\pm|_U\|_2\le\sqrt{2(1+\eta)d}.
\]
For this bound, use \(|g(h_\pm)|\le\|P_Ah_\pm\|_2\) and factor
the difference of squares.  Its other factor is at most \(2(1+\eta)\).
The coefficients satisfying \(\alpha_\pm(t)g(b_t)\le0\) have total
mass at most \(d\).

Retain only the other coefficients ending in \(U\), and call their
intrinsic forest contributions \(v_\pm\). Then
\[
 y|_U=v_++w_+,\qquad -y|_U=v_-+w_-,\qquad
 L_U(w_\pm)\le K:=\sqrt{2(1+\eta)d}+d.
\]
Within each tree of \(U\), the value \(g(b_t)\) is constant, since
the head of every such path is the same. If that value is zero, no
coefficient is retained there. Otherwise both retained arrays have its
sign. One isometric change of signs makes both
arrays nonnegative. Lemma~\ref{lem:linear-positive-monotonicity} now gives
\[
 L_U(v_+)\le L_U(v_++v_-)
 =L_U(w_++w_-)\le2K.
\]
It follows that \(L_U(y|_U)\le3K\) and \(L_n(y)\le6K\).
Letting \(\eta\downarrow0\), and using
\(0\le1-r\le\sqrt{1-r}\), proves
\eqref{eq:linear-positive-unit}. Scaling gives the first bound in
\eqref{eq:linear-positive-scaled} when \(\beta>0\); when \(\beta=0\),
both vectors are zero. Contractivity of \(P_A\) gives \(r\le\beta\)
and the second bound. In each component both endpoint norms are at
least the center norm. Thus the maximum of their squared differences
from \(r^2\) is at most their sum. Summing over components proves
\eqref{eq:linear-positive-outer}, including components with empty heads.
\end{proof}

\subsection{Removing mass by subtree quotas}
\label{sec:linear-quota-removal}

There is another way to use the consistent signs outside a head. Instead
of changing a norming functional, remove a small portion of the path
coefficients. The remaining path sum can then be bounded in Hilbert norm.
The following finite selection fact is the common combinatorial step.

\begin{lemma}[Subtree quotas]
\label{lem:linear-subtree-quota}
On a forest of finite height, let \(\mu,\nu\) be nonnegative finitely
supported endpoint masses. For a vertex \(v\), let \(K_v\) be its
descendant set, and put \(G(v)=\mu(K_v)\), \(W(v)=\nu(K_v)\).
There is \(0\le\rho\le\mu\) such that
\[
 \rho(K_v)\ge\min\{G(v),W(v)\}\quad\hbox{for every }v,
 \qquad \|\rho\|_1\le\|\nu\|_1.
\]
The last inequality also holds separately in each tree of the forest.
\end{lemma}
\begin{proof}
Only the finite ancestor hull of the supports matters. Write
\(q(v)=\min\{G(v),W(v)\}\). Nonnegativity gives
\[
 q(v)\le G(v),\qquad
 \sum_{w\text{ child of }v}q(w)\le q(v).
\]
Process vertices from the leaves towards the roots. At \(v\), the mass
already marked in its child subtrees totals \(\sum_wq(w)\). Mark enough
additional, previously unmarked mass of \(\mu\) below \(v\) to make
the marked total there equal to \(q(v)\); fractions of an endpoint mass
are allowed. The two displayed inequalities make this possible. Later
steps can only increase a previously met quota. The total finally marked
in a tree equals its root quota and is at most its total \(\nu\)-mass.
Take \(\rho\) to be the marked mass.
\end{proof}

\begin{theorem}[The bottom-up tail estimate]
\label{thm:linear-bottom-up-tail}
Suppose \(0\le\varepsilon\le1\), \(P_Ax=x\), \(P_Ay=0\),
\(L_n(x\pm y)\le1\), and \(L_n(x)\ge1-\varepsilon\). Then
\begin{equation}\label{eq:linear-bottom-up-explicit}
 L_n(y)\le3\sqrt{2\varepsilon}+8\varepsilon.
\end{equation}
Moreover the following uniform implication holds on \(X\): for any
sequence of finite head projections \(P_j\) and vectors satisfying
\[
 P_jx_j=x_j,\quad P_jy_j=0,\quad
 \|x_j\pm y_j\|\le1,\quad\|x_j\|\longrightarrow1,
\]
one has \(\|y_j\|\longrightarrow0\).
\end{theorem}
\begin{proof}
Choose \(g=gP_A\) norming \(x\), and finite path representations of
both endpoints with cost at most \(1+\eta\). Put
\(d=1+\eta-L_n(x)\le\varepsilon+\eta\). As in the preceding proof,
their Hilbert tails have lengths \(K_\pm\le\sqrt{2(1+\eta)d}\), and
their coefficient masses \(B_\pm\) with
\(\alpha_\pm(t)g(b_t)\le0\) satisfy \(B_\pm\le d\). Call these
coefficients bad and the others good.

In a complementary tree where the common value \(c=g(b_t)\) is zero,
there are no good coefficients. Otherwise put \(\varsigma=\operatorname{sign}(c)\).
All good coefficients have sign \(\varsigma\), and all bad nonzero
coefficients have sign \(-\varsigma\). Apply
Lemma~\ref{lem:linear-subtree-quota} to the combined good mass \(\mu\)
and combined bad mass \(\nu\) in each tree, and mark the selected
portions with their originating endpoint representation. Denote the
selected good mass by \(\rho\), and put
\(G(v)=\mu(K_v)\), \(W(v)=\nu(K_v)\). At a vertex in a tree with \(c\ne0\),
the vanishing sum of the two endpoint tails reads
\[
 h_+(v)+h_-(v)+\varsigma\bigl(G(v)-W(v)\bigr)=0.
\]
Consequently, in every complementary tree, including those with no good mass,
\[
 (G(v)-W(v))_+\le|h_+(v)+h_-(v)|\qquad(v\notin A).
\]
After removal, the good tail in either one of the representations has
absolute coordinate at most \(G(v)-\rho(K_v)\le(G(v)-W(v))_+\).
Its Hilbert norm is therefore at most \(K_++K_-\). The total removed
mass is at most \(B_++B_-\), and each complementary path costs at most
two in \(L_n\). Decomposing the plus endpoint tail into its Hilbert,
bad, removed-good, and remaining-good parts gives
\[
 L_n(y)\le2K_++K_-+2B_++2(B_++B_-)
 \le3\sqrt{2(1+\eta)d}+8d.
\]
Let \(\eta\downarrow0\). This proves
\eqref{eq:linear-bottom-up-explicit}, including its endpoint cases.

We give the outer-sum passage explicitly. For the \(j\)-th pair set
\begin{align*}
 a_{j,n}^\pm&=L_n(x_{j,n}\pm y_{j,n}),
 &r_{j,n}&=L_n(x_{j,n}),\\
 c_{j,n}&=\tfrac12(a_{j,n}^++a_{j,n}^-),
 &a_{j,n}&=\max(a_{j,n}^+,a_{j,n}^-).
\end{align*}
Convexity gives \(c_{j,n}\ge r_{j,n}\), and the endpoint assumptions
give \(\|a_j^\pm\|_2\le1\). Hence \(\|c_j\|_2\to1\). The Hilbert
parallelogram identity then gives
\[
 \|a_j^+-a_j^-\|_2^2
 =2\|a_j^+\|_2^2+2\|a_j^-\|_2^2-4\|c_j\|_2^2
 \le4(1-\|x_j\|^2)\longrightarrow0.
\]
Since \(a_j=c_j+|a_j^+-a_j^-|/2\), we also have \(\|a_j\|_2\to1\).
Ignore coordinates with \(a_{j,n}=0\), whose vectors are zero, and put
\(\delta_{j,n}=1-r_{j,n}/a_{j,n}\in[0,1]\). Then
\[
 \sum_na_{j,n}^2\delta_{j,n}
 \le\sum_n(a_{j,n}^2-r_{j,n}^2)\longrightarrow0.
\]
Applying the component estimate after division by \(a_{j,n}\) gives
\(L_n(y_{j,n})/a_{j,n}\le3\sqrt{2\delta_{j,n}}+8\delta_{j,n}\);
the same ratio is at most one by the triangle inequality. For fixed
\(\tau>0\), the squared outer norm is therefore at most
\[
 (3\sqrt{2\tau}+8\tau)^2\sum_na_{j,n}^2
 +\tau^{-1}\sum_na_{j,n}^2\delta_{j,n}.
\]
First let \(j\to\infty\), then \(\tau\downarrow0\).
\end{proof}

The preceding argument emphasizes uniform decay. A different threshold
in the same selection procedure gives a convenient homogeneous squared
estimate, without an outer limiting argument.

\begin{theorem}[A squared estimate by quota removal]
\label{thm:linear-trimming-242}
For a finite head projection \(P\) on \(X\),
\begin{equation}\label{eq:linear-trimming-242}
 \|y\|^2\le242\bigl(\|x+y\|^2+\|x-y\|^2-2\|x\|^2\bigr)
 \qquad(Px=x,\ Py=0).
\end{equation}
The same assertion, with \(L_n\) in place of \(\|\cdot\|\), holds in
each component.
\end{theorem}
\begin{proof}
Work first in one component. Set \(r=L_n(x)\) and
\(M_\pm=L_n(x\pm y)\), so \(M_\pm\ge r\). Use a supported norming
functional \(g\) and finite representations of costs at most
\(M_\pm+\eta\). Put \(d_\pm=M_\pm+\eta-r\). Here a path coefficient
is good when \(\operatorname{sign}(\alpha_\pm(t))g(b_t)>1/2\), and bad
otherwise. The deficit calculation gives
\begin{equation}\label{eq:linear-quota-deficits}
 B_\pm\le2d_\pm,\qquad
 K_\pm\le\sqrt{2(M_\pm+\eta)d_\pm}.
\end{equation}
Write \(c=g(b_t)\) for the common value on one complementary tree.
If \(|c|\le1/2\), that tree has no good coefficients. Otherwise put
\(\varsigma=\operatorname{sign}(c)\); good coefficients have sign
\(\varsigma\), and bad nonzero coefficients have sign \(-\varsigma\).
Let \(G_+(v),G_-(v)\) be the subtree masses of the good coefficients
from the two representations, and let \(W(v)\) be the subtree mass of
all bad coefficients. In a tree with \(|c|>1/2\), the endpoint tails satisfy
\[
 h_+(v)+h_-(v)+\varsigma\bigl(G_+(v)+G_-(v)-W(v)\bigr)=0.
\]
For the quota selection take \(G=G_+\). Since \(G_-\ge0\), the identity
gives the following bound, which is also immediate in a tree with no good mass:
\[
 (G(v)-W(v))_+\le|h_+(v)+h_-(v)|.
\]
Lemma~\ref{lem:linear-subtree-quota} removes at most \(B_++B_-\)
from the plus good coefficients. The remaining plus good tail has
Hilbert norm at most \(K_++K_-\); the removed tail has \(L_n\)-norm
at most \(2(B_++B_-)\). Adding the plus Hilbert and bad tails yields
\[
 \begin{split}
 L_n(y)&\le2(K_++K_-)+4(B_++B_-)\\
 &\le11\sum_{\sigma\in\{+,-\}}
             \sqrt{(M_\sigma+\eta)d_\sigma}.
 \end{split}
\]
Here \(d_\sigma\le M_\sigma+\eta\) and
\(2\sqrt2+8<11\). Squaring the last sum costs at most a factor two,
and
\((M_\sigma+\eta)d_\sigma\le(M_\sigma+\eta)^2-r^2\).
Letting \(\eta\downarrow0\) proves the component estimate with
\(2\cdot11^2=242\). Summation gives
\eqref{eq:linear-trimming-242}. Empty heads cause no exception; their
centers are zero.
\end{proof}

\subsection{Stopping at the first excess of wrong-sign mass}
\label{sec:linear-stopped-fibers}

For a displacement supported in a larger finite head, its part in the
smaller head need not vanish. We now average the endpoint decompositions before choosing signs, and
measure the mass disagreeing with this average in each head fibre.
Stopping where that mass exceeds one quarter of the total gives an
estimate for the new coordinates alone.

\begin{theorem}[A fresh-layer squared estimate]
\label{thm:linear-fresh-layer-400}
Let \(A\subseteq A'\) be finite ancestor-closed subsets of the disjoint
union of the \(T_n\), and put \(E=A'\setminus A\). If
\(P_Ax=x\) and \(P_{A'}y=y\), then
\begin{equation}\label{eq:linear-fresh-layer-400}
 \|P_Ey\|^2\le400\left(
       \frac{\|x+y\|^2+\|x-y\|^2}{2}-\|x\|^2\right).
\end{equation}
There is no condition on \(P_Ay\).
\end{theorem}
\begin{proof}
Write \(x^s=x+sy\), \(s\in\{-1,1\}\), and denote averages over
\(s\) by a bar. Fix \(\lambda>0\). In each nonzero component choose
\[
 x_n^s=z_n^s+B_n\alpha_n^s,\qquad
 r_n^s=\|z_n^s\|_2+M_n^s\le(1+\lambda)L_n(x_n^s),\qquad
 M_n^s=\|\alpha_n^s\|_1.
\]
Use zero decompositions in zero components. Project to \(A'\): truncate
the Hilbert terms and push each path endpoint to its last ancestor in
\(A'\), discarding endpoints with no such ancestor. Both costs decrease,
so all arrays may be taken supported in \(A'\).

In a fixed \(T_n\), let \(\pi_A(t)\) be the last ancestor of \(t\)
in \(A\), or a dummy symbol \(\perp\) if none exists. For each fibre set
\[
 L_b=\sum_{\pi_A(t)=b}\bar\alpha_n(t),\qquad
 l_b\in\{-1,1\},\quad l_bL_b=|L_b|.
\]
Either choice is allowed at zero. Introduce the mass on pairs \((s,t)\)
given by \(m_n(s,t)=|\alpha_n^s(t)|/2\), and let \(h(s,t)\) be the
coefficient's sign, chosen arbitrarily at zero. Its wrong-sign mass is
\[
 W_n=\sum_{s,t:\,h(s,t)\ne l_{\pi_A(t)}}m_n(s,t).
\]
Separating the two signs in each fibre, including the dummy fibre, gives
\(\bar M_n-\sum_b|L_b|=2W_n\). Since
\(x_n=P_{A\cap T_n}(\bar z_n+B_n\bar\alpha_n)\), its component norm is at most
\(\overline{\|P_Az_n^s\|_2}+\sum_b|L_b|\); the dummy term is a
harmless nonnegative addition to this cost. It follows that
\begin{equation}\label{eq:linear-stopped-deficit}
 \Delta_n:=\bar r_n-L_n(x_n)
 \ge2W_n+\frac12\overline{
             \frac{\|P_Ez_n^s\|_2^2}{\|z_n^s\|_2}}\ge0,
\end{equation}
A quotient with zero denominator means zero.  The Hilbert estimate
used here is
\[
 \|z\|_2-\|P_Az\|_2\ge\frac{\|P_Ez\|_2^2}{2\|z\|_2}
 \quad(z\ne0,\ P_{A'}z=z).
\]

The quantity
\(Q=\sum_n\bar r_n\Delta_n\) controls both Hilbert leakage and
wrong-sign mass. For an array on the disjoint union of the trees,
\(\|\cdot\|_2\) denotes its Hilbert norm, including summation over
components. Since \(\|z_n^s\|_2\le2\bar r_n\),
\begin{equation}\label{eq:linear-stopped-Q}
 \begin{split}
 \|P_E\bar z\|_2^2
 &\le\overline{\|P_Ez^s\|_2^2}\le4Q,\qquad
 \sum_n\Delta_n^2\le Q,\\
 Q&\le\sum_n(\bar r_n)^2-\|x\|^2
 \le(1+\lambda)^2\overline{\|x^s\|^2}-\|x\|^2.
 \end{split}
\end{equation}
The first line follows from \eqref{eq:linear-stopped-deficit} and
\(0\le\Delta_n\le\bar r_n\); the second uses
\(L_n(x_n)\le\bar r_n\) and
\((\bar r_n)^2\le\overline{(r_n^s)^2}\).

It remains to recover the norm of \(P_Ey\) from this budget. The zero
case is immediate. Otherwise take a norming functional for \(P_Ey\)
and precompose it with \(P_E=P_{A'}-P_A\). Its coefficient array
\(\phi\), supported in \(E\), satisfies
\begin{equation}\label{eq:linear-stopped-dual}
 \langle y,\phi\rangle=\|P_Ey\|,\qquad
 \|\phi\|_2\le2,\qquad
 \Bigl(\sum_nc_n^2\Bigr)^{1/2}\le2,
 \quad c_n=\sup_{t\in T_n}\left|\sum_{v\preceq t}\phi(v)\right|.
\end{equation}
The functional has norm at most two. Testing against Hilbert vectors
gives the first coefficient bound. Testing against finite sums of one
root path in each component, with arbitrary scalar multipliers, gives
the second by the outer Hilbert norm. Approximate each supremum if it
is not attained.

The Hilbert part of
\(\langle y,\phi\rangle=\overline{s\langle x^s,\phi\rangle}\)
has absolute value at most \(4\sqrt Q\), by
\eqref{eq:linear-stopped-Q} and \eqref{eq:linear-stopped-dual}.
For the path part write \(g(t)=\sum_{v\preceq t}\phi(v)\) and
\[
 I_n=\int s\,h(s,t)\,g(t)\,dm_n(s,t).
\]
For \(v\in E\cap T_n\), let \(m(v)\) be the mass of pairs with
\(t\succeq v\), and \(w(v)\) its wrong-sign part. The entire subtree
below \(v\) is in a single \(\pi_A\)-fibre. Its coefficients therefore satisfy
\[
 \bigl(B_n\bar\alpha_n\bigr)(v)
 =l_{\pi_A(v)}\bigl(m(v)-2w(v)\bigr)=-\bar z(v),
\]
where the last equality uses \(x(v)=0\) on \(E\). Thus, whenever
\(w(v)\le m(v)/4\),
\[
 m(v)\le2|\bigl(B_n\bar\alpha_n\bigr)(v)|=2|\bar z(v)|.
\]

Call a vertex \(v\in E\) retained if this inequality's condition holds
at \(v\) and at every preceding vertex of its path in \(E\). Split
\(g(t)\) into its sum over retained vertices and the remainder. The
remainder starts at the first vertex satisfying \(w(v)>m(v)/4\).
Such first vertices form an antichain, so their disjoint descendant
sets have total mass at most \(4W_n\). The remainder is the difference
of two prefix sums and has absolute value at most \(2c_n\). It therefore
contributes at most \(8c_nW_n\) to \(|I_n|\). The retained contribution
is at most
\(2\sum_{v\in E\cap T_n}|\phi(v)|\,|\bar z(v)|\).
Summing over components and using \(2W_n\le\Delta_n\), we obtain
\[
 \begin{split}
 \sum_n|I_n|
 &\le4\|P_E\bar z\|_2+4\sum_nc_n\Delta_n\\
 &\le8\sqrt Q+8\sqrt Q=16\sqrt Q.
 \end{split}
\]
Together with the Hilbert part this proves \(\|P_Ey\|\le20\sqrt Q\).
Apply \eqref{eq:linear-stopped-Q}, square, and let
\(\lambda\downarrow0\) to obtain \eqref{eq:linear-fresh-layer-400}.
\end{proof}

\section{Energy on newly exposed coordinates}\label{sec:linear-energy}

Throughout this section \(X\) is the outer Hilbert sum in
\eqref{eq:linear-outer-sum}, with block norm \(L_n\) on
\(H_n=\ell_2(T_n)\), root-path vectors \(b_u\), and path map \(B_n\).
Proposition~\ref{prop:linear-duality} identifies \(L_n\) with the
supremum over the tests
\begin{equation}\label{eq:linear-energy-tests}
 \|a\|_2\le1,\qquad |\langle b_u,a\rangle|\le1\quad(u\in T_n).
\end{equation}
All coordinate sets below lie in the disjoint union of the trees.  A set
is \emph{initial} if it contains the ancestors of each of its vertices.
We write \(R_A=P_A\) for coordinate restriction to \(A\).  Recall that
\(R_A\) is contractive for initial \(A\), and that
\begin{equation}
 L_n(R_{B\setminus A}b_u)\le2\qquad(A\subseteq B\text{ initial}).
 \label{eq:linear-window-path}
\end{equation}
Indeed, a root path meets \(B\setminus A\) in an interval, whose
indicator is the difference of two root-path indicators, with a missing
endpoint interpreted as zero.

The estimates in this section measure how much of a dyadic martingale
can appear on successive new layers.  Predictability means that the
layer is chosen before the sign producing the new martingale value.
It is this timing, together with the coherence of path coefficients,
that makes the estimates uniform in the tree heights.  We first prove
a bound for martingales supported on predictable heads.  We then give a
second proof using variation, and two estimates that allow leakage of
the parent vector into the new window. The first proof establishes a
signed comparison for its child and parent mixtures, using the same
positive-mass idea as above in a direct polar argument. The variation
proof applies the subtree-quota lemma to its collapsed coefficients.

\subsection{Predicting the coefficient of one random path}

Fix an integer \(m\ge1\).  On \(\Omega=\{-1,1\}^m\), with uniform
probability, let \(\mathcal F_i\)
be the sigma-field generated by the first \(i\) signs, and write
\(\mathbb E_i\) for conditional expectation.  An \(X\)-valued dyadic
martingale is a sequence \(x_i=\mathbb E_i x_m\).

\begin{theorem}[Predictable-layer energy]\label{thm:linear-predictable-energy}
Let \(x_0,\ldots,x_m\) be an \(X\)-valued dyadic martingale.  Let \(P_0\)
be a fixed finite initial set, and, for \(1\le i\le m\), let \(P_i\)
be a finite initial set measurable with respect to \(\mathcal F_{i-1}\).
Suppose, on every history, that
\[
 P_{i-1}\subseteq P_i\quad(1\le i\le m),\qquad
 \operatorname{supp}x_i\subseteq P_i\quad(0\le i\le m),\qquad \|x_m\|\le C.
\]
Then, for \(\Delta_i=P_i\setminus P_{i-1}\),
\begin{equation}\label{eq:linear-energy-495}
 \sum_{i=1}^m\mathbb E\|R_{\Delta_i}x_i\|^2\le495C^2.
\end{equation}
\end{theorem}

\begin{proof}
There are only finitely many relevant coordinates, since there are
finitely many sign histories and all the sets \(P_i\) are finite.
We first represent the terminal value by one random atom in each
component.  The finite-head unit ball is exactly the convex hull in
\eqref{eq:linear-finite-ball}.  This description provides finite random
representations by a Hilbert-ball vector or a signed root path.

For each terminal history and each component, apply this description
with \(P=P_m\cap T_n\), rescaled by \(L_n(x_m^n)\).  On a finite
extension of the probability space we may therefore choose \(h_n,c_n,u_n\)
so that
\begin{equation}\label{eq:linear-one-atom}
 \mathbb E[h_n+c_nb_{u_n}\mid\mathcal F_m]=x_m^n,
 \qquad
 \sum_n(\|h_n\|_2^2+|c_n|^2)\le C^2.
\end{equation}
For each component the sample is either a Hilbert-ball vector
\(h_n\), with \(c_n=0\), or a signed root path with \(h_n=0\) and
\(|c_n|=L_n(x_m^n)\).  Zero components use both coefficients zero.
An arbitrary root can be used for \(u_n\) when \(c_n=0\).  The
sigma-fields \(\mathcal F_i\) still record only the original signs.
In particular \(x_i=\mathbb E_i(h_n+c_nb_{u_n})_n\).

Let \(v_i^n\) be the deepest vertex of the path to \(u_n\) in \(P_i\),
or a symbol \(\partial\) if this intersection is empty; set
\(b_\partial=0\).  The path trace on the new layer is
\[
 I_i^n=R_{\Delta_i\cap T_n}b_{u_n}=b_{v_i^n}-b_{v_{i-1}^n},
 \qquad L_n(I_i^n)\le2.
\]
The sigma-fields \(\mathcal G_i^n=\sigma(\mathcal F_i,v_i^n)\) increase:
from \(v_i^n\) and the previous sign history one recovers \(v_{i-1}^n\)
by truncating its path to \(P_{i-1}\).  Set
\(a_i^n=\mathbb E[c_n\mid\mathcal G_i^n]\).  Orthogonality of scalar
martingale increments gives
\begin{equation}\label{eq:linear-coefficient-energy}
 \sum_n\sum_{i=1}^m\mathbb E|a_i^n-a_{i-1}^n|^2
 \le\sum_n\mathbb E|c_n|^2\le C^2.
\end{equation}

We split the new layer into its Hilbert contribution, its old predicted
path coefficient, and the change in that prediction:
\begin{align*}
 R_{\Delta_i}x_i&=H_i+M_i+W_i,\\
 H_i&=\mathbb E_iR_{\Delta_i}(h_n)_n,\\
 M_i&=\mathbb E_i(a_{i-1}^nI_i^n)_n,\\
 W_i&=\mathbb E_i((a_i^n-a_{i-1}^n)I_i^n)_n.
\end{align*}
The identity follows because \(I_i^n\) is \(\mathcal G_i^n\)-measurable,
so \(c_n\) may be replaced by \(a_i^n\) inside its conditional
expectation.  Conditional Jensen, the disjointness of the \(\Delta_i\) on
each outcome, and \eqref{eq:linear-coefficient-energy} show that
\begin{equation}\label{eq:linear-HW-energy}
 \sum_i\mathbb E\|H_i\|^2\le C^2,
 \qquad \sum_i\mathbb E\|W_i\|^2\le4C^2.
\end{equation}
Here we use \(L_n\le\|\cdot\|_2\) for the first inequality.  Denote
conditional expectations given \(\mathcal F_{i-1}\) by primes.
Predictability of \(\Delta_i\), together with the support of \(x_{i-1}\),
gives \(H_i'+M_i'+W_i'=R_{\Delta_i}x_{i-1}=0\).  Consequently
\begin{equation}\label{eq:linear-old-M-energy}
 \sum_i\mathbb E\|M_i'\|^2\le10C^2.
\end{equation}
It remains to compare the old-coefficient contribution on a child
history with its average on the parent history.

We prove \(\|M_i\|\le4\|M_i'\|\) pointwise.  Fix a component and a
history through time \(i-1\).  For \(v\in \Delta_i\cap T_n\), let \(r(v)\)
be the first vertex of \(\Delta_i\) on its root path.  If \(v_i^n=v\), then
\(v_{i-1}^n\) is the parent of \(r(v)\), or \(\partial\) when this
first vertex is the root. Thus the value of \(a_{i-1}^n\) on these
outcomes is a scalar \(q_{r(v)}\) depending only on this first vertex.
It is the conditional mean on the old trace cell determined by the
parent history and \(v_{i-1}^n=r(v)^-\), with the root case interpreted
as \(\partial\). On a null old trace cell put \(q_{r(v)}=0\); all its
parent and child trace probabilities then vanish.
Let \(\mu(v)\) and \(\widetilde\mu(v)\) be the probabilities of
\(v_i^n=v\) conditional on this parent history and on a specified child
history, respectively.  Since both children have probability \(1/2\),
\(\widetilde\mu(v)\le2\mu(v)\).

Put \(I(v)=b_v-b_{r(v)^-}\), interpreting the second term as zero when
\(r(v)\) is the root. Outcomes with \(v_i^n\notin\Delta_i\) have zero
path trace. The two conditional mixtures are therefore
\[
 \begin{aligned}
 M_i^n&=\sum_{v\in\Delta_i\cap T_n}q_{r(v)}\widetilde\mu(v)I(v),\\
 (M_i')^n&=\sum_{v\in\Delta_i\cap T_n}q_{r(v)}\mu(v)I(v).
 \end{aligned}
\]

For an admissible test \(a\), put
\(A_v=\sum_{w\in[r(v),v]}a_w\).  We have
\[
 |\langle M_i^n,a\rangle|
 \le2\sum_{v\in \Delta_i\cap T_n}|q_{r(v)}|\mu(v)|A_v|.
\]
Define a test candidate \(f\), zero off \(\Delta_i\cap T_n\), by
\[
 f_v=\operatorname{sign}(q_{r(v)})(|A_v|-|A_{v^-}|),
\]
where \(A_{v^-}=0\) at the first vertex of each interval and
\(\operatorname{sign}(0)=0\).  The reverse triangle inequality gives
\(\|f\|_2\le\|a\|_2\).  Each root path meets \(\Delta_i\) in at most one
interval.  Summing \(f\) on that interval telescopes to a signed
\(|A_v|\), which is at most \(2\), since \(A_v\) is the difference of
two root-path sums of \(a\).  Hence \(f/2\) is admissible.  Moreover
\[
 \langle(M_i')^n,f\rangle
 =\sum_{v\in \Delta_i\cap T_n}|q_{r(v)}|\mu(v)|A_v|.
\]
Taking the supremum over \(a\) proves
\(L_n(M_i^n)\le4L_n((M_i')^n)\).  Summing the component squares proves
the required comparison.  By \eqref{eq:linear-old-M-energy}, the total
energy of \(M_i\) is at most \(160C^2\).  Finally,
\(\|H_i+M_i+W_i\|^2\le3(\|H_i\|^2+\|M_i\|^2+\|W_i\|^2)\), giving
\(3(1+160+4)C^2=495C^2\).
\end{proof}

\subsection{Variation and removal of positive mass}

A different argument keeps all terminal path coefficients rather than
sampling one path.  The quantity that telescopes is then the variation
of their collapsed coefficients.  Its gain controls the amount of
positive mass that must be removed to make a child layer small in the
Hilbert norm.

For a finite initial \(R\subset T_n\), define
\[
 d_R(z,u)=z(u)-\sum_{\substack{v\text{ child of }u\\v\in R}}z(v),
 \qquad V_R(z)=\sum_{u\in R}|d_R(z,u)|,
 \qquad V_\varnothing(z)=0.
\]
Finite summation up the tree gives
\begin{equation}\label{eq:linear-variation-representation}
 R_Rz=\sum_{u\in R}d_R(z,u)b_u,
 \qquad V_R(B_nc)\le\|c\|_1.
\end{equation}
For the second inequality, each collapsed coefficient sums the \(c_v\)
whose deepest ancestor in \(R\) is its label; these classes are disjoint.

\begin{lemma}[Removing excess mass]\label{lem:linear-mass-removal}
Let \(S\subseteq T\subset T_n\) be finite initial sets, let
\(z^+,z^-\in H_n\), and put \(z=(z^++z^-)/2\).  Put \(U=T\setminus S\) and
\[
 g=\tfrac12\big(V_T(z^+)+V_T(z^-)\big)-V_S(z).
\]
Then \(g\ge0\) and, for either sign,
\begin{equation}\label{eq:linear-mass-removal}
 L_n(R_Uz^\pm)\le2\|R_Uz\|_2+4g.
\end{equation}
\end{lemma}

\begin{proof}
If \(S=\varnothing\), then
\(L_n(R_Tz^\pm)\le V_T(z^\pm)\le2g\), so the conclusion follows.
Otherwise \(S\) contains the root.  For \(u\in T\), write \(\pi(u)\)
for its deepest ancestor in \(S\).  Choose
\(\sigma_s\in\{-1,1\}\) so that
\(\sigma_sd_S(z,s)=|d_S(z,s)|\).  Write
\[
 d_T(z^\pm,u)=\sigma_{\pi(u)}(p_u^\pm-q_u^\pm),
 \qquad p_u^\pm,q_u^\pm\ge0,
\]
using the positive and negative parts.  Collapsing the refined path
coefficients to \(S\) shows that
\begin{equation}\label{eq:linear-variation-gap}
 g=2\sum_{u\in T}\tfrac12(q_u^++q_u^-)\ge0.
\end{equation}
For \(u\in U\), all descendants in \(T\) have the same \(\pi\)-label.
Let \(P_u^\pm,Q_u^\pm\) be the respective sums of \(p_v^\pm,q_v^\pm\)
over these descendants.  Since
\(z^\pm(u)=\sigma_{\pi(u)}(P_u^\pm-Q_u^\pm)\),
\begin{equation}\label{eq:linear-positive-domination}
 P_u^\pm\le2\operatorname{av}P_u^\pm
 \le2|z(u)|+2\operatorname{av}Q_u^\pm,
\end{equation}
where \(\operatorname{av}\) means the average of the two signs.

Fix one sign and apply Lemma~\ref{lem:linear-subtree-quota} on the finite
forest \(U\) to the endpoint masses \(\mu_u=p_u^\pm\) and
\(\nu_u=q_u^++q_u^-=2\operatorname{av}q_u^\pm\). Initiality of \(S\)
ensures that all descendants in \(T\) of a vertex in \(U\) remain in
\(U\). The corresponding subtree masses are therefore
\(P_u^\pm\) and \(2\operatorname{av}Q_u^\pm\). The lemma gives
\(0\le r_u\le p_u^\pm\) with
\[
 \sum_{v\succeq u}r_v
 \ge\min\{P_u^\pm,2\operatorname{av}Q_u^\pm\},
 \qquad
 \sum_{u\in U}r_u\le2\sum_{u\in U}\operatorname{av}q_u^\pm\le g.
\]
By \eqref{eq:linear-positive-domination}, the remaining mass satisfies
\[
 P_u^\pm-\sum_{v\succeq u}r_v
 \le\bigl(P_u^\pm-2\operatorname{av}Q_u^\pm\bigr)_+
 \le2|z(u)|.
\]
Also \(\sum_{u\in U}q_u^\pm\le g\), by
\eqref{eq:linear-variation-gap}.  The remaining positive mass defines
\[
 v(u)=\sigma_{\pi(u)}\sum_{w\succeq u}(p_w^\pm-r_w)
 \quad(u\in U),\qquad v=0\text{ off }U.
\]
Its coordinates satisfy \(|v(u)|\le2|z(u)|\).  The exact decomposition
\[
 R_Uz^\pm=v+\sum_{u\in U}\sigma_{\pi(u)}(r_u-q_u^\pm)R_Ub_u
\]
therefore costs at most \(2\|R_Uz\|_2+2(g+g)\), as required.
\end{proof}

\begin{theorem}[Variation energy]\label{thm:linear-variation-energy}
Under all hypotheses of
Theorem~\ref{thm:linear-predictable-energy},
\begin{equation}\label{eq:linear-energy-252}
 \sum_{i=1}^m\mathbb E\|R_{\Delta_i}x_i\|^2\le252C^2.
\end{equation}
In particular the bound is \(1008\) when \(\|x_m\|\le2\).
\end{theorem}

\begin{proof}
For each terminal history choose, componentwise,
\[
 x_m^n=h_n+B_nc_n,
 \qquad \|h_n\|_2+\|c_n\|_1\le2L_n(x_m^n),
\]
with both terms zero in a zero component.  Only finitely many
components are needed.  Put \(h_i=\mathbb E_i(h_n)_n\),
\(b_i=(B_n\mathbb E_i c_n)_n\); then \(x_i=h_i+b_i\).  The terminal
Hilbert and coefficient budgets are each at most \(4C^2\):
\[
 \mathbb E\|(h_n)_n\|_{\ell_2(\coprod T_n)}^2\le4C^2,
 \qquad \mathbb E\sum_n\|c_n\|_1^2\le4C^2.
\]
At time \(i-1\), apply Lemma~\ref{lem:linear-mass-removal} in each
component to \(S=P_{i-1}\cap T_n\), \(T=P_i\cap T_n\), and the two
values of \(b_i^n\).  These sets are fixed before the sign at time
\(i\).  If \(g_{i,n}\) is the variation gain in the lemma, then
\[
 V_{P_{i-1}\cap T_n}(b_{i-1}^n)^2+g_{i,n}^2
 \le\mathbb E_{i-1}V_{P_i\cap T_n}(b_i^n)^2.
\]
Indeed, the average variation equals the previous variation plus the
nonnegative number \(g_{i,n}\); square and apply scalar Jensen.
Telescoping and \eqref{eq:linear-variation-representation} give
\[
 \sum_{i,n}\mathbb E g_{i,n}^2\le4C^2.
\]
Since \(R_{\Delta_i}x_{i-1}=0\), the lemma and the outer Euclidean norm imply
\[
 \|R_{\Delta_i}x_i\|
 \le\|R_{\Delta_i}h_i\|_2+2\|R_{\Delta_i}h_{i-1}\|_2
      +4\Big(\sum_ng_{i,n}^2\Big)^{1/2}.
\]
For \(j=i-1\) and \(j=i\), predictability gives
\(R_{\Delta_i}h_j=\mathbb E_jR_{\Delta_i}(h_n)_n\).
Conditional Jensen followed by pathwise disjointness of the layers
bounds each of the two unscaled Hilbert energies by \(4C^2\).
Squaring the last three-term inequality and summing now gives
\(3(1+4+16)4C^2=252C^2\).
\end{proof}

\subsection{Outcome-dependent path measures and sign prediction}

The next estimate permits an error on the parent window.  Here both the
positive path mass and its sign may depend on the entire outcome.  We
place them on one measure space of outcome--path pairs.  The scalar
martingale of conditional signs then pays for changes of sign
prediction, while a stopping antichain controls the remaining positive
path mass.

\begin{theorem}[Adaptive-window energy]\label{thm:linear-sign-energy}
There is an absolute constant \(K\) with the following property.
Let \(x_t=\mathbb E_t x_m\), \(0\le t\le m\), be an \(X\)-valued dyadic
martingale.  For \(0\le t<m\), let \(A_t\subseteq D_t\) be finite
initial sets measurable with respect to \(\mathcal F_t\), and suppose
\(D_t\subseteq A_{t+1}\) whenever \(t+1<m\).  Put
\(U_t=D_t\setminus A_t\).  Then
\begin{equation}\label{eq:linear-sign-energy}
 \sum_{t=0}^{m-1}\mathbb E\|R_{U_t}x_{t+1}\|^2
 \le K\left(\mathbb E\|x_m\|^2+
                  \sum_{t=0}^{m-1}\mathbb E\|R_{U_t}x_t\|^2\right).
\end{equation}
No support condition on \(x_t\) is imposed.
\end{theorem}

\begin{proof}
For every terminal outcome and component choose
\[
 x_m^n=h_n+B_n\lambda_n,
 \qquad \|h_n\|_2+\|\lambda_n\|_1\le2L_n(x_m^n).
\]
Write \(\lambda_n=\epsilon_n\mu_n\), where \(\mu_n=|\lambda_n|\),
\(|\epsilon_n|\le1\), and \(M_n=\mu_n(T_n)\).  Thus \(\mu_n\) is a
positive measure on path endpoints, allowed to depend on the terminal
outcome.  The two budgets satisfy
\begin{equation}\label{eq:linear-sign-budgets}
 \mathbb E\sum_n(\|h_n\|_2^2+M_n^2)\le4\mathbb E\|x_m\|^2.
\end{equation}
On outcome--endpoint pairs put the finite positive measure
\[
 \nu_n(\{\omega\}\times\{u\})=2^{-m}\mu_n(\omega,\{u\}).
\]
This need not be a probability measure; conditional means on its
positive-mass cells have their usual meaning. On null cells set them
to zero. The operators \(\mathbb E_t\) still use the original uniform
law on sign histories, whereas the conditional means of \(\epsilon_n\)
below use \(\nu_n\), which need not give equal mass to the two sign
children. Let \(a_{t,n}^-\) be the conditional mean of \(\epsilon_n\)
given the first \(t\) signs and the truncation of the root path to
\(A_t\).  Let \(a_{t,n}^+\) be its conditional mean given the first
\(t+1\) signs and its truncation to \(D_t\).  Empty truncations are
recorded by a dummy symbol.  The conditioning partitions, ordered as
\[
 (0,-),(0,+),(1,-),(1,+),\ldots,(m-1,-),(m-1,+),
\]
refine successively: the known sign history and truncation of a deeper
trace to the earlier head recover the preceding data. Define
\[
 q_{t,n}=\mathbb E_{t+1}
       \int|a_{t,n}^+-a_{t,n}^-|\,d\mu_n.
\]
We first establish the weighted sign estimate
\begin{equation}\label{eq:linear-weighted-sign}
 \sum_{t<m}\mathbb E q_{t,n}^2\le4\mathbb E M_n^2.
\end{equation}

For clarity fix \(n\) temporarily.  Conditional Cauchy--Schwarz gives
\[
 q_t^2\le(\mathbb E_{t+1}M)\,
       \mathbb E_{t+1}\int|a_t^+-a_t^-|^2\,d\mu.
\]
Under the original sign law the children are fair, so \(M\ge0\) gives
\(\mathbb E_{t+1}M\le2\mathbb E_tM\).  Put
\[
 R_t=\max_{0\le j\le t}\mathbb E_jM,\qquad R_{-1}=0.
\]
The conversion from the sign law to the pair measure is
\[
 \mathbb E\left[(\mathbb E_tM)\,
   \mathbb E_{t+1}\int|a_t^+-a_t^-|^2\,d\mu\right]
 =\int(\mathbb E_tM)|a_t^+-a_t^-|^2\,d\nu,
\]
because \(\mathbb E_tM\) is measurable at time \(t+1\). Hence the
sum of the integrated right sides is at most
\[
 2\sum_{t<m}\int R_t|a_t^+-a_t^-|^2\,d\nu.
\]
Expand \(R_t=\sum_{j\le t}(R_j-R_{j-1})\). Each nonnegative increment
is constant on each cell \(C\) of the partition defining \(a_j^-\).
Orthogonality of successive conditional-mean increments under \(\nu\)
gives
\[
 \begin{aligned}
 \sum_{t=j}^{m-1}\int_C|a_t^+-a_t^-|^2\,d\nu
 &\le\int_C|\epsilon-a_j^-|^2\,d\nu\\
 &=\int_C|\epsilon|^2\,d\nu-\nu(C)|a_j^-|^2
 \le\nu(C).
 \end{aligned}
\]
The left side uses only some increments of the refining sequence;
omitting the others preserves the bound. Multiply by the cell-constant
value of \(R_j-R_{j-1}\), then sum over cells and \(j\). The weighted
sum \(2\sum_{t<m}\int R_t|a_t^+-a_t^-|^2\,d\nu\) is consequently at most
\(2\int R_{m-1}\,d\nu=2\mathbb E(MR_{m-1})\).

The scalar input is the finite-horizon form of Doob's $L_2$ maximal
inequality~\cite{Doob1953}.  We include its first-crossing proof:

\[
 \Big\|\max_{j\le m}\mathbb E_jM\Big\|_2\le2\|M\|_2.
\]
Indeed, if \(R=\max_{j\le m}\mathbb E_jM\), partitioning the event
\(\{R\ge u\}\) according to its first crossing gives
\(u\mathbb P(R\ge u)\le\mathbb E[M\mathbf1_{\{R\ge u\}}]\).
Integrating in \(u\) yields
\(\mathbb ER^2\le2\mathbb E(MR)\); Cauchy--Schwarz proves the claim.
Thus \(2\mathbb E(MR_{m-1})\le4\mathbb EM^2\), proving
\eqref{eq:linear-weighted-sign}.

We next compare a child window with its parent window, keeping the
component \(n\) fixed.  For a vertex \(s\), write
\(C_s=\{u:s\preceq u\}\) for the cone of endpoints whose paths pass
through \(s\). Define the signed endpoint measures
\[
 \lambda'_{t+1}=\mathbb E_{t+1}(a_t^-\mu_n),
 \qquad
 \rho_{t+1}=\mathbb E_{t+1}((a_t^+-a_t^-)\mu_n).
\]
For \(s\in U_t\), the child sign history together with membership
in \(C_s\) is measurable for the partition defining \(a_t^+\). Hence
\[
 \mathbb E_{t+1}\int_{C_s}\epsilon_n\,d\mu_n
 =\mathbb E_{t+1}\int_{C_s}a_t^+\,d\mu_n
 =\lambda'_{t+1}(C_s)+\rho_{t+1}(C_s).
\]
Thus their sum represents the path part of \(x_{t+1}^n\) on the window
\(U_t\). The error measure satisfies \(\|\rho_{t+1}\|_1\le q_{t,n}\).
Let \(h_{t,n}=\mathbb E_th_n\),
\(\lambda'_t=\mathbb E_t\lambda'_{t+1}\), and
\(\rho_t=\mathbb E_t\rho_{t+1}\).  Averaging gives, at each
\(s\in U_t\),
\begin{equation}\label{eq:linear-parent-path}
 x_t^n(s)=h_{t,n}(s)+\lambda'_t(C_s)+\rho_t(C_s).
\end{equation}

Fix a history through time \(t\).  On every cone \(C_s\) with
\(s\in U_t\), truncation to \(A_t\) is fixed.  Thus \(a_t^-\) is
constant on that cone, including all completions of the fixed sign
history.  There is no cancellation within the cone in either
\(\lambda'_t\) or \(\lambda'_{t+1}\).  With
\(\phi_s=|\lambda'_t(C_s)|\), fair binary conditioning therefore gives
\begin{equation}\label{eq:linear-positive-path-comparison}
 |\lambda'_{t+1}|(C_s)\le2\phi_s.
\end{equation}
Choose another path decomposition, now only for the parent window:
\[
 R_{U_t}x_t^n=d+B_n\theta,
 \qquad \|d\|_2+\|\theta\|_1\le2r_{t,n},
 \qquad r_{t,n}=L_n(R_{U_t}x_t^n).
\]
If \(\beta=|\theta|+|\rho_t|\), then
\eqref{eq:linear-parent-path} gives
\begin{equation}\label{eq:linear-stopping-domination}
 \phi_s\le|d(s)-h_{t,n}(s)|+\beta(C_s),
 \qquad \|\beta\|_1\le\|\theta\|_1+\mathbb E_tq_{t,n}.
\end{equation}

Stop along each descending line at the first vertex \(s\in U_t\)
where \(\beta(C_s)\ge\phi_s/2\).  These first vertices form an
antichain, so their endpoint cones are disjoint.  At window vertices
neither stopped nor below a stop,
\(\phi_s\le2|d(s)-h_{t,n}(s)|\).  By
\eqref{eq:linear-positive-path-comparison}, the vector formed by the
child's path coordinates at these vertices has Hilbert norm at most
\(4(\|d\|_2+\|R_{U_t}h_{t,n}\|_2)\).
On the remaining part of the window, every contributing endpoint lies
in a stopping cone.  Its total \(|\lambda'_{t+1}|\)-mass is at most
\(4\|\beta\|_1\), using the stopping inequality, the factor \(2\) in
\eqref{eq:linear-positive-path-comparison}, and disjointness of the
cones.  Each path contributes there on an interval, of \(L_n\)-norm
at most \(2\).  This part therefore costs at most \(8\|\beta\|_1\).
The error measure \(\rho_{t+1}\) costs at most \(2q_{t,n}\).
Combining these estimates gives the explicit component bound
\begin{equation}\label{eq:linear-sign-window}
 \begin{split}
 L_n(R_{U_t}x_{t+1}^n)\le{}&
 \|R_{U_t}h_{t+1,n}\|_2+4\|R_{U_t}h_{t,n}\|_2\\
 &+16r_{t,n}+8\mathbb E_tq_{t,n}+2q_{t,n}.
 \end{split}
\end{equation}
Here \(4\|d\|_2+8\|\theta\|_1\le16r_{t,n}\).

We have reduced the vector estimate to square summation.  Conditional
Jensen and predictability of \(U_t\) give, for \(j=t,t+1\),
\[
 \mathbb E\|R_{U_t}h_{j,n}\|_2^2
 \le\mathbb E\|R_{U_t}h_n\|_2^2.
\]
The windows are disjoint on every outcome, so summing over \(t,n\)
bounds either Hilbert energy by \(\mathbb E\sum_n\|h_n\|_2^2\).
The same summation for \(q_{t,n}\) and \(\mathbb E_tq_{t,n}\) is
bounded by \(4\mathbb E\sum_nM_n^2\), by
\eqref{eq:linear-weighted-sign} and Jensen.  Square the five terms in
\eqref{eq:linear-sign-window}, bounding their sum's square by five
times the sum of squares, and use \eqref{eq:linear-sign-budgets}.
This proves \eqref{eq:linear-sign-energy} with an absolute constant.
All sums over components involve nonnegative terms; Tonelli's theorem
justifies them even when the terminal vectors have infinitely many
nonzero components.
\end{proof}

\subsection{Fixed windows and a direct leakage estimate}

For fixed heads, a simpler collapse of the terminal signed masses gives
an explicit leakage constant.  It uses the mass whose sign disagrees
with the previous collapsed coefficient, rather than a martingale of
predicted signs.  Unlike Theorem~\ref{thm:linear-predictable-energy},
the conclusion pays for the actual parent contribution to each window.

\begin{theorem}[Energy with leakage]\label{thm:linear-leakage-energy}
Let \(W_i=\mathbb E_iW_m\), \(0\le i\le m\), be an \(X\)-valued dyadic
martingale.  Let \(A_0\subseteq\cdots\subseteq A_m\) be fixed finite
initial sets and put \(\Delta_i=A_i\setminus A_{i-1}\).  Then
\begin{equation}\label{eq:linear-leakage-1024}
 \sum_{i=1}^m\mathbb E\|R_{\Delta_i}W_i\|^2
 \le1024\left(\mathbb E\|W_m\|^2+
              \sum_{i=1}^m\mathbb E\|R_{\Delta_i}W_{i-1}\|^2\right).
\end{equation}
\end{theorem}

\begin{proof}
We first work in one component, suppressing \(n\) and restricting each
set to its tree.  At each terminal outcome choose
\[
 W_m=H+B_nc,
 \qquad \|H\|_2+\|c\|_1\le2L_n(W_m),
\]
and let \(H_i=\mathbb E_iH\).  For an endpoint \(s\), let \(\pi_i(s)\)
be its deepest ancestor in \(A_i\), or \(\partial\) if none exists.
Set \(\pi_i(\partial)=\partial\) and \(b_\partial=0\).  Define
\[
 c_r^i=\mathbb E_i\sum_{\pi_i(s)=r}c_s,
 \qquad v_i=\sum_{r\in A_i\cup\{\partial\}}|c_r^i|.
\]
The dummy coefficient is kept in this last sum even though its path
vector vanishes.  This ensures the exact refinement identity
\begin{equation}\label{eq:linear-collapsed-refinement}
 \mathbb E_{i-1}\sum_{\pi_{i-1}(r)=q}c_r^i=c_q^{i-1}
 \quad(q\in A_{i-1}\cup\{\partial\}).
\end{equation}
The child window itself has the representation
\[
 R_{\Delta_i}W_i=R_{\Delta_i}H_i+\sum_rc_r^iR_{\Delta_i}b_r.
\]

Orient each refined coefficient by its old collapsed coefficient:
\[
 \sigma_r=\operatorname{sign}(c^{i-1}_{\pi_{i-1}(r)}),
 \qquad c_r^i=\sigma_r(g_r-q_r),\qquad g_r,q_r\ge0,
\]
where sign at zero is chosen to be \(+1\), and \(g_r,q_r\) are the
positive and negative parts of \(\sigma_rc_r^i\).  The signs are known
on the parent history.  By \eqref{eq:linear-collapsed-refinement},
\[
 \beta_i:=\mathbb E_{i-1}\sum_rq_r
       =\tfrac12(\mathbb E_{i-1}v_i-v_{i-1})\ge0.
\]
The elementary inequality
\((\mathbb E_{i-1}v_i-v_{i-1})^2
 \le\mathbb E_{i-1}v_i^2-v_{i-1}^2\)
therefore telescopes to
\begin{equation}\label{eq:linear-bad-mass-energy}
 \sum_i\mathbb E\beta_i^2\le\tfrac14\mathbb Ev_m^2
          \le\tfrac14\mathbb E\|c\|_1^2.
\end{equation}

Fix a step and a parent history.  For \(t\in \Delta_i\), all its descendant
labels in \(A_i\) have the same old ancestor and therefore the same
sign \(\sigma_t\).  Set
\(g(t)=\sum_{r\succeq t}g_r\) and
\(q(t)=\sum_{r\succeq t}q_r\), with the dummy omitted.  Averaging the
coordinate formula gives
\[
 W_{i-1}(t)=H_{i-1}(t)+
     \sigma_t(\mathbb E_{i-1}g(t)-\mathbb E_{i-1}q(t)).
\]
Decompose the leaked parent window once more:
\[
 Y=R_{\Delta_i}W_{i-1}=L+B_nd,
 \qquad \|L\|_2+\|d\|_1\le2L_n(Y).
\]
There are two equally likely children and \(g(t)\ge0\), so
\[
 g(t)\le2\mathbb E_{i-1}g(t)
 \le2\left(|L(t)|+|H_{i-1}(t)|+
                 \sum_{s\succeq t}|d_s|+\mathbb E_{i-1}q(t)\right).
\]
Call \(t\in \Delta_i\) large when
\(g(t)>4(|L(t)|+|H_{i-1}(t)|)\).  At a large node,
\[
 g(t)\le4\left(\sum_{s\succeq t}|d_s|+
                         \mathbb E_{i-1}q(t)\right).
\]
The minimal large nodes form an antichain.  Mark every coefficient
label in \(A_i\) below one of them.  Their descendant sets are disjoint,
so the total marked positive mass is at most
\(4(\|d\|_1+\beta_i)\).  The total negative mass satisfies
\(\sum_rq_r\le2\beta_i\), again by fair binary averaging.
Using \(L_n(R_{\Delta_i}b_r)\le2\), all marked positive terms together with
all negative terms cost at most \(8\|d\|_1+12\beta_i\).

It remains to estimate the unmarked positive part.  Its coordinate is
zero at every large vertex, since all descendant labels there are
marked.  At other vertices its absolute value is at most
\(g(t)\le4(|L(t)|+|H_{i-1}(t)|)\).  Its Hilbert norm is consequently
at most \(4(\|L\|_2+\|R_{\Delta_i}H_{i-1}\|_2)\).  Including the child
Hilbert term yields
\begin{equation}\label{eq:linear-leakage-one-step}
 \begin{split}
 L_n(R_{\Delta_i}W_i)\le{}&\|R_{\Delta_i}H_i\|_2
     +4\|R_{\Delta_i}H_{i-1}\|_2+12\beta_i\\
     &+16L_n(R_{\Delta_i}W_{i-1}).
 \end{split}
\end{equation}
The last coefficient follows from
\(4\|L\|_2+8\|d\|_1\le16L_n(Y)\).

For \(j=i-1,i\), conditional Jensen and the fixed, disjoint windows give
\(\sum_i\mathbb E\|R_{\Delta_i}H_j\|_2^2\le\mathbb E\|H\|_2^2\).
Squaring the four terms in \eqref{eq:linear-leakage-one-step}, using
four times their squared sum, and applying
\eqref{eq:linear-bad-mass-energy}, gives
\[
 \sum_i\mathbb E L_n(R_{\Delta_i}W_i)^2
 \le68\mathbb E\|H\|_2^2+144\mathbb E\|c\|_1^2
        +1024\sum_i\mathbb E L_n(R_{\Delta_i}W_{i-1})^2.
\]
The first two terms are at most
\(144\mathbb E(\|H\|_2+\|c\|_1)^2
 \le576\mathbb E L_n(W_m)^2\).
Finally sum over the components and replace \(576\) by \(1024\) in the
terminal term.  This proves \eqref{eq:linear-leakage-1024}.
\end{proof}
